% !PRINTED 697
\weberpage{697}
%<<CONT
of this remark one may, if $n : \mu_n$ approaches a finite limiting value $\gamma$, take $\alpha$ and $\beta$ as near to this limiting value as one pleases, and so obtains the somewhat more precise formulation of the theorem:

3. \gesperrt{If the positive numbers $\mu_1, \mu_2, \mu_3, \dots$ are such that}
\[
\mathop{\mathrm{Lim}}\limits_{n=\infty} \frac{n}{\mu_n} = \gamma
\]
\gesperrt{is a finite limiting value, then the infinite series}
\[
S = \frac{1}{\mu_1^s} + \frac{1}{\mu_2^s} + \frac{1}{\mu_3^s} + \cdots
\]
\gesperrt{is convergent for every $s$ that is greater than 1, and we have}
\[
\mathop{\mathrm{Lim}}\limits_{s=1} (s-1)\,S = \gamma.
\]

Let there now be given any system $M$ of \gesperrt{infinitely many} positive numbers $\mu_n$, which we order thus
\weq{5}{\mu_1 \leqq \mu_2 \leqq \mu_3 \leqq \dots \qquad (M),}
so that in the sequence (5) a smaller term never follows a greater one, and let $T$ denote the number of terms in $M$ that are not greater than an arbitrarily given positive quantity $t$. The following theorem then holds:

4. \gesperrt{If one of the two limiting values}
\weq{6}{\mathop{\mathrm{Lim}}\limits_{n=\infty} \frac{n}{\mu_n}, \quad \mathop{\mathrm{Lim}}\limits_{t=\infty} \frac{T}{t}}
\gesperrt{is finite, then the other has the same finite value.}

Let first
\weq{7}{\mathop{\mathrm{Lim}}\limits_{n=\infty} \frac{n}{\mu_n} = \gamma}
be a finite limiting value. Then the $\mu_n$ must necessarily grow to infinity as $n$ increases without limit, and for every positive $t$ there is a value $m$ such that
\weq{8}{\mu_m \leqq t < \mu_{m+1};}
$m$ increases to infinity together with $t$, and one has $T = m$ [because of (5)]. Hence by (8)
\weq{9}{\frac{m}{\mu_m} \geqq \frac{T}{t} > \frac{m}{\mu_{m+1}}.}
